\chapter{Graph}

\section{Fundamentals}
	\kactlimport{BellmanFord.h}
	% \kactlimport{FloydWarshall.h}
	% \kactlimport{TopoSort.h}

\section{Network flow}
	\subsection{Edge lower bounds}
		Say we require $L(e) \leq f(e) \leq R(e)$ for every edge $e$. If $R(e) < 0$, reverse the edge, and flip the signs. If $L(e) < 0 \le R(e)$, set $L(e)$ back to $0$, but add a reverse edge with capacity $-L(e)$. Add super source and sink $s'$ and $t'$. Replace edge $a \to b$ with three edges:
			\begin{itemize}
				\item Same edge with capacity $R(e) - L(e)$.
				\item Add edge $s' \to b$ with capacity $L(e)$.
				\item Add edge $a \to t'$ with capacity $L(e)$.
			\end{itemize}
		Add (original sink to source) $t\to s$ with infinite capacity. Feasible iff $(s', t')$ has max flow equal to $\sum L(e)$. (((Minimal flow of original network: binary search on the capactiy of $t\to s$ (instead of infinity). This edge captures the max allowed total flow of the original network.))) After finishing feasibility run, run max flow again \textbf{over} the remainder of the modified graph, from original source to sink for original max flow value. Add $L(e)$ to each edge for the exact flow of each edge. Use Dinic or MCMF instead of PushRelabel as they allow re-running.
		% \vspace{1px}
	\kactlimport{PushRelabel.h}
	\kactlimport{MinCostMaxFlow.h}
	% \kactlimport{EdmondsKarp.h}
	\kactlimport{Dinic.h}
	\kactlimport{MinCut.h}
	\kactlimport{GlobalMinCut.h}
	% \kactlimport{GomoryHu.h}

\section{Matching}
	\kactlimport{hopcroftKarp.h}
	\kactlimport{DFSMatching.h}
	\kactlimport{MinimumVertexCover.h}
	\kactlimport{WeightedMatching.h}
	\kactlimport{GeneralMatching.h}

\section{DFS algorithms}
	\kactlimport{SCC.h}
	\kactlimport{BiconnectedComponents.h}
	\kactlimport{2sat.h}
	\kactlimport{EulerWalk.h}

\section{Coloring}
	\kactlimport{EdgeColoring.h}
	\kactlimport{Sack.h}

\section{Heuristics}
	\kactlimport{MaximalCliques.h}
	\kactlimport{MaximumClique.h}
	\kactlimport{MaximumIndependentSet.h}

\section{Trees}
	% \kactlimport{BinaryLifting.h}
	\kactlimport{LCA.h}
	\kactlimport{CompressTree.h}
	\kactlimport{CentroidDecomp.h}
	\kactlimport{HLD.h}
	\kactlimport{LinkCutTree.h}
	\kactlimport{DirectedMST.h}

\newpage

\section{Math}
	\subsection{Number of Spanning Trees}
		% I.e. matrix-tree theorem.
		% Source: https://en.wikipedia.org/wiki/Kirchhoff%27s_theorem
		% Test: stress-tests/graph/matrix-tree.cpp
		Create an $N\times N$ matrix \texttt{mat}, and for each edge $a \rightarrow b \in G$, do
		\texttt{mat[a][b]--, mat[b][b]++} (and \texttt{mat[b][a]--, mat[a][a]++} if $G$ is undirected).
		Remove the $i$th row and column and take the determinant; this yields the number of directed spanning trees rooted at $i$
		(if $G$ is undirected, remove any row/column).
		% \vspace{25px}

	\subsection{Erdős–Gallai theorem}
		% Source: https://en.wikipedia.org/wiki/Erd%C5%91s%E2%80%93Gallai_theorem
		% Test: stress-tests/graph/erdos-gallai.cpp
		A simple graph with node degrees $d_1 \ge \dots \ge d_n$ exists iff $d_1 + \dots + d_n$ is even and for every $k = 1\dots n$, $\sum _{i=1}^{k}d_{i}\leq k(k-1)+\sum _{i=k+1}^{n}\min(d_{i},k)$
